\documentclass[12 pt]{elsarticle}

\usepackage{caption}
\usepackage{cancel}
\usepackage{amsmath}
\usepackage{amssymb}
\usepackage{amsfonts}
\usepackage{graphicx}
\usepackage[T1]{fontenc}
\usepackage[utf8]{inputenc}
\usepackage{rotating}
\usepackage{booktabs}
%\usepackage{graphics}
\usepackage[margin=1.3in]{geometry}
\usepackage{subfig}
\usepackage{setspace}
\usepackage{appendix}
\usepackage{threeparttable}
\usepackage{hyperref}
\hypersetup{colorlinks=true,
      linkcolor=red,
      citecolor=black,
      filecolor=magenta,
      urlcolor=cyan,
      pdftitle={},
      pdfauthor={},
      pdfsubject={},
      pdfkeywords={},
      }
\usepackage{color}
\usepackage{lscape}
\usepackage{umoline}
\usepackage{multirow}
\usepackage{fancyhdr}
%\usepackage{subcaption}
\usepackage{chngcntr}
\usepackage{etoolbox}
\usepackage{lipsum}
\usepackage{caption}
\usepackage{pdfpages}
\usepackage{chronosys}
\usepackage{tikz}
\usetikzlibrary{positioning,arrows.meta}
\usepackage{epigraph}
\usetikzlibrary{decorations.pathreplacing}
\makeatletter
\def\ps@pprintTitle{
 \let\@oddhead\@empty
 \let\@evenhead\@empty
 \def\@oddfoot{\reset@font\hfil\thepage\hfil}
 \let\@evenfoot\@oddfoot
}
\makeatother
\setlength{\bibsep}{0pt plus 0.3ex}
\doublespacing
\biboptions{authoryear}
\setcounter{MaxMatrixCols}{10}

\usepackage{titlesec}
\titleformat*{\section}{\Large\bfseries}
\titleformat*{\subsection}{\large\bfseries}
\titleformat*{\subsubsection}{\large\bfseries}
\titleformat*{\subparagraph}{\large\bfseries}

\linespread{1.3}
\newtheorem{theorem}{Theorem}
\newtheorem{acknowledgement}[theorem]{Acknowledgement}
\newtheorem{algorithm}[theorem]{Algorithm}
\newtheorem{axiom}[theorem]{Axiom}
\newtheorem{case}[theorem]{Case}
\newtheorem{claim}[theorem]{Claim}
\newtheorem{conclusion}[theorem]{Conclusion}
\newtheorem{condition}[theorem]{Condition}
\newtheorem{conjecture}[theorem]{Conjecture}
\newtheorem{corollary}[theorem]{Corollary}
\newtheorem{criterion}[theorem]{Criterion}
\newtheorem{definition}[theorem]{Definition}
\newtheorem{example}[theorem]{Example}
\newtheorem{exercise}[theorem]{Exercise}
\newtheorem{lemma}[theorem]{Lemma}
\newtheorem{notation}[theorem]{Notation}
\newtheorem{problem}[theorem]{Problem}
\newtheorem{proposition}[theorem]{Proposition}
\newtheorem{remark}[theorem]{Remark}
\newtheorem{solution}[theorem]{Solution}
\newtheorem{summary}[theorem]{Summary}
\newenvironment{proof}[1][Proof]{\textbf{#1.} }{\ \rule{0.5em}{0.5em}}
\renewcommand {\baselinestretch} {1.3}
\geometry{left=28mm,right=30mm,top=28mm,bottom=28mm}
\interfootnotelinepenalty=50000
\newcommand\cites[1]{\citeauthor{#1}'s\ (\citeyear{#1})}
\renewcommand\theaffn{\arabic{affn}}
\usepackage{lscape} 

\begin{document}
\begin{frontmatter}
\title{\textbf{Beneath the Surface: Investigating Anti-Competitive Signals in Brazilian Public Procurement\tnoteref{tfoot}}}\tnotetext[tfoot]{
\emph{This version: July 2024 (First version: July 2024)}}
\author[ETH]{Sergio Galletta}
\ead{sergio.galletta@gess.ethz.ch}
\author[INSPER]{Darcio Genicolo-Martins}
\ead{darciogm1@insper.edu.br}
\author[ETH]{Tommaso Giommoni}
\ead{giommoni@ethz.ch}
\address[ETH]{ETH Zürich}
\address[INSPER]{INSPER}


\begin{abstract}
In this study, we adopt a Regression Discontinuity design, following the methodology of \cite{kawai2023using}, to investigate 3.8 million public procurement auctions in São Paulo, Brazil. We aim to detect bid rotation and incumbency trends as indicators of anti-competitive behavior by analyzing auctions with narrowly winning and narrowly losing bids. Our findings reveal an incumbency advantage but no definitive evidence of bid rotation. This indicates that political favoritism or corruption, rather than firm collusion, may influence the auction outcomes.

%\newline
\begin{center}
-- Very Preliminary and Incomplete --   \\
-- Please, Do Not Circulate --   
\end{center}
\end{abstract}
\end{frontmatter}

\bigskip

\clearpage

%citation to add: buhler2022state, dehdari2022origins, couttenier2019violent, jha2012does, dell2018nation, gehring2022can, becker2020wars, okunogbe2022technology, cage2020heroes 

\section{Introduction}\label{intro}
The integrity and efficiency of procurement processes are crucial for ensuring the effective use of public funds, supporting economic development, and maintaining public trust \citep{Bosio2022}. %\footnote{Public procurement accounts for about 12\% of GDP worldwide and up to 20\% in developing countries, making it a vital component of national economies \citep{Bosio2022}.} 
Anti-competitive behavior in public procurement can increase the costs for governments and taxpayers and reduce the quality and affordability of public services \citep{Sokol2017}. It also disrupts fair competition, limits innovation, and puts honest businesses at a disadvantage. 
Detecting anti-competitive behavior in procurement is difficult due to its hidden nature and the complexity of the processes involved. Ineffective procurement can result from several factors, including corruption and political favoritism, which can influence contract awards. Additionally, private companies may form cartels to raise prices and reduce competition. While most research on procurement in developing countries focuses on corruption, less attention is given to detecting collusive behavior.\footnote{Only, \cite{kawai2022study} Indonesia, \cite{baranek2021detection} and ukrainia}

In this paper, we use the methodology suggested by \cite{kawai2023using} and employ a Regression Discontinuity design to detect patterns indicative of anti-competitive behavior in São Paulo's public procurement data, covering 3.8 million contracts from 2009 to 2019. Our research aims to identify instances where firms winning by narrow margins display unusual incumbency advantages, possibly due to political favoritism, or have low backlogs, which may indicate bid rotation schemes among competitors.


Our findings reveal that firms winning by a small margin typically have a stronger track record of past competition victories. Specifically, we find that winning firms have a history of winning more often than losing firms. We also find that losing firms tend to have significantly lower backlogs, with their backlogs being XX\% lower on average than winning firms. 

Our paper contributes to the existing evidence by highlighting anti-competitive conditions in public procurement within the Brazilian context. It suggests that these issues are likely not only due to corruption and political favoritism but also involve more conventional collusion cartels by firms. 




%Collusive activities are usually hidden under layers of corporate secrecy and embedded deep within market structures, making them difficult to detect without advanced analytical tools \citep{MarshallMarx2012}.




%In Brazil, particularly in the state of São Paulo, the economic center of the country, significant financial resources are allocated each year for public procurement. This underscores the importance of detecting and preventing collusion in these auctions for economic efficiency and ethical governance.


%The state of São Paulo is particularly vulnerable to the risks of procurement collusion due to the large amount and scale of its procurement activities. Recent improvements in Brazil's anti-corruption laws and the implementation of stricter procurement regulations highlight the urgent need to address these issues \citep{OECD2016}.


The paper proceeds as follows. Section \ref{coll} describes our methods. \ref{hist} and \ref{data} provide institutional details and information on the data source, respectively. Sections \ref{res} and \ref{res_ar} present our results. Section \ref{concl} concludes. 

 \section{Identifying anti-competitive behaviors (TBC)}\label{coll}

To detect potential anticompetitive behaviors in public procurements, we adopt the methodology proposed by \cite{kawai2023using}. This strategy is based on the assumption that companies submitting nearly identical bids should, under competitive conditions, have similar cost structures and historical success rates. These closely matched bids suggest that the winning probabilities should be almost identical unless external factors such as collusion or political favoritism are at play.

Practically, we implement a Regression Discontinuity Design (RDD), where the key running variable is the difference between the winning bid and the closest losing bid, denoted as $\Delta_{i,t}$, with a cutoff set at zero. This difference sets a natural threshold: bids marginally lower than the highest losing bid win the auction, whereas those slightly higher do not. The main variables of interest in our analysis are the incumbency status and the backlog of the bidding firms,  which should reveal any anomalies if non-competitive practices influence the outcomes.

We hypothesize that if non-competitive practices, such as preferential treatment or collusive arrangements, influence auction outcomes, this would be evident through significant discontinuities in the incumbency rates and backlogs of firms that bid around the cutoff point. For example, in non-competitive environments, due to favoritism, we should expect unusually high past success rates for firms with bids just below the cutoff compared to those just above it.

Formally, the coefficient \(\beta\) quantifies this expected discontinuity, defined as:
\[
\beta = \lim_{\epsilon \to 0^+} E[x_{i,t} | \Delta_{i,t} = \epsilon] - \lim_{\epsilon \to 0^-} E[x_{i,t} | \Delta_{i,t} = \epsilon].
\]
Here, \(\Delta_{i,t} = b_{i,t} - \land b_{-i,t}\) indicates the difference between the bid of firm \(i\) and the most competitive alternative bid at time \(t\). If \(\Delta_{i,t} < 0\), then bidder \(i\) is the winner of the auction; conversely if \(\Delta_{i,t} > 0\), then bidder \(i\) does not win. The variable \(x_{i,t}\) measures our outcome of interest of a firm \(i\)'s up until time \(t\).

\section{Institutional background (TBC)} \label{hist}

This section provides a brief institutional background on public tenders in the state of São Paulo, Brazil, focusing on details relevant to our empirical analysis.

Each individual government entity responsible for procurement, the Public Buyer Unit (PBU), purchases goods and services in a decentralized manner through the Bolsa Eletrônica de Compras (BEC), the state's e-procurement platform. Since 2007, it has been mandatory to use the BEC for purchasing common goods and services, including for all 99 SES/SP units. The BEC plays a significant role in SES/SP procurement. In 2018 alone, approximately USD 1.7 billion were traded through the platform. Since 2009, the BEC has facilitated more than BRL 7.4 billion in SES/SP negotiations, accounting for approximately 34\% of all state purchases. The planning and execution of tenders consist of a very costly public administration process that requires relevant financial and human resources. An acquisition can have administrative costs from US\$500 to US\$5,200, depending on the bid complexity. Thus, if a public tender does not obtain a supplier, it creates relevant waste for the government.


%This section provides a brief institutional background on public tenders in the state of Sao Paulo, Brazil. We focus on the details that are most relevant for the empirical analysis.

%Each Public Buyer Unit (PBU) purchases in a decentralized way through the Bolsa Eletronica de Compras (BEC), the e-procurement platform of São Paulo state. Since 2007, it has been mandatory to use the BEC to purchase common goods and services in São Paulo state, including all 99 SES/SP units. The BEC figures of SES/SP buying are expressive. In 2018, approximately US\$1.7 billion were traded and since 2009, the electronic platform has handled more than R\$7.4 billion in SES/SP negotiations (approximately 34\% of all state purchases).

Nearly 55,000 item purchase offers were negotiated, covering approximately 6,350 different items. The BEC features an extensive catalog of standardized items and services, all described in great detail. The SES/SP utilizes the same two types of competitive tendering procedures available on the BEC for both ordinary and urgent purchases: (i) sealed bid tendering (convite) and (ii) multi round descending auctions (pregão).

%Almost 55,000 item purchase offers were negotiated, comprising approximately 6,350 different traded items. The BEC has an extensive catalog of standardized items and services that are described in great detail. The SES/SP uses the same two types of competitive tendering procedures available in the BEC to make ordinary and urgent purchases (litigated and administrative): (i) sealed bid tendering (convite) and (ii) multiround descending auctions (pregão).


In sealed bid tenders, firms submit their proposals to the government by a specified date mentioned in the notice. Once the submission period ends, the proposals and participant details are made public when the auctioneer opens the envelopes. The contract is awarded to the firm with the appropriate documentation that submitted the lowest bid. The sealed bid tendering process, is allowed for purchases up to a value limit of R\$176,000 (approximately US\$35,200).


%In sealed bids, firms send their proposals to the government on a specific date specified in the notice. Later, the proposals and participants become public information when the auctioneer 'opens' the envelopes. The winning firm is the one among those with the appropriate documentation that submitted the lowest bid. Convite is allowed up to the purchase value limit of R\$176,000 (approximately US\$35,200).

The multi round descending auction method has no limit on the purchase value and combines elements of a modified sealed-bid tender and a reverse auction. Initially, PBUs evaluate and classify qualified proposals based on the conditions outlined in the sealed bid phase notice. The auctioneer then publicly reveals all valid proposals while keeping the firms' identities anonymous. Following this, the descending auction begins: for 20 minutes, each qualified firm submits bids, aware of the current lowest valid bid. If a valid bid is made between 16 and 20 minutes, the auction is extended by 4 minutes. The auction concludes only if 4 minutes pass without a new valid bid. The final criterion for awarding the tender is the lowest bid price, which must be below the reference price.


%Pregão has no limit on the purchase value. This mode is a combination of a modified sealed-bid tender and reverse auction. In this case, PBUs classify qualified proposals only under the conditions set out in the sealed bid phase notice. Then, the auctioneer publicly reveals all valid proposals, keeping the firms' identities anonymous. Next, the descending auction begins: for 20 minutes, each qualified firm submits its bids, knowing the current lowest valid bid. If there is a valid bid between 16 and 20 minutes, the auction will be extended for another 4 minutes. It ends only if 4 minutes pass without a valid bid. The final criterion for winning the tender is presenting the lowest bid price, which must be lower than the reference price.


In sum, the main difference between the sealed bid tender and the multi-round descending auction is that the latter allows for a negotiation phase after the reverse auction, during which companies and the government can negotiate the lowest final price previously obtained. In contrast, the sealed bid tender has a single phase, making it a more straightforward procedure to perform and monitor.

%One of the main differences between the convite and pregão modes is that the latter allows for a negotiation phase after the reverse auction during which companies and the government can negotiate the lowest final price previously obtained. On the other hand, since convite has a single phase, it tends to be a more straightforward procedure to perform and monitor.




\section{Data (TBC)}\label{data}
\subsection{Sources} We use administrative data on bid-level public procurement tenders of common goods and services in the state of São Paulo, Brazil, from January 2009 to December 2019. These transactions were carried out on the BEC electronic procurement platform, which is available to all PBUs throughout the state. The Department of Finance of São Paulo state (SEFAZ/SP) is responsible for the operational management and centralization of the BEC bidding data.

In total, 1,344 PBUs make regular purchases through the BEC, including state-level executive, legislative, and judiciary bureaus in the state of São Paulo as well as other affiliated entities, such as some municipalities located in the state of São Paulo and a group of private organizations. PBUs purchased 169,607 different types of items (goods and services), totaling 3,866,407 transactions from 19,007 different firms in this period.

The BEC has a very detailed catalog of standardized goods and services organized in three levels: group, class, and item. For instance, health items are classified as group 65 (Medical, dental, and hospital equipment and supplies). Thus, the item coded 110639 is the drug 'Furosemide 40 milligrams, coated tablets, units' belonging to class 6531 (Medicines prescribed with or without ANVISA notification/registration) and group 65.

Data are organized by purchase offer (PO), the electronic document issued by the PBU that identifies and quantifies the goods and services that will be purchased. A PO is defined by a 22-character code and may contain one or more listed items, but each item has its own purchase process. Thus, the purchase of an item is uniquely identified by the combination of the PO and the purchased item codes (POI).

Each POI provides information on the internal phase parameters, such as item quantities and reference prices, and external phase outcomes, such as bid prices (winners and losers), the number of participating firms, the number of bids, whether the public tender was successful or not, and the identification of the PBU and the auctioneer, among other variables.
\subsection{Variables definition}
We define our variables of interest following the methodology outlined by \cite{kawai2022study}. The first variable is a measure of \textit{incumbency}, represented by a dummy variable that takes 1 if a firm won the last auction it participated in and 0 otherwise. The second variable measures \textit{backlog}. To construct this measure, we first sum each firm's total amount of awarded work. Then, following \cite{kawai2022study}, we standardize this raw backlog within each firm using its mean and standard deviation. This standardization accounts for intertemporal changes and firm size differences, ensuring a more accurate representation of the backlog.
Finally, we define the \textit{margin of victory} as the percentage difference between the lowest and second-lowest bids, calculated as a percentage of the lowest bid. This variable is used as the running variable in the RDD setting.


%\section{Empirical Strategy}\label{emp}


%We want to explore the effects of government audits on anti-competitive behavior in procurement. To do so, we leverage the natural experiment provided by Brazil's random audits program. This program randomly selects municipalities to audit. Abundant research has been generated based on this government program, but no evidence about the actual effect on bidding behavior has been provided.

%We argue that audits create an exogenous boost in the perceived likelihood of detection of misconduct. Theoretically, the increase in the perceived probability of being caught due to these audits should deter anti-competitive practices, such as collusion and overpricing in public procurements.

%Our empirical strategy involves comparing municipalities that were audited with those that were not, using the random nature of audit selection to control for unobservable factors that could influence misconduct. By focusing on changes in bidding behavior and contract awards before and after audits, we aim to identify any shifts that could indicate a reduction in anti-competitive practices. 

%To assess the impact of government audits on anti-competitive behavior we use a Difference-in-Differences (DiD) approach formalized as follow:

%\begin{equation}
 %   Y_{it} = \alpha + \beta_1 \text{Audit}_{it} + \beta_2 \text{Post}_{t} + \delta (\text{Audit}_{it} \times \text{Post}_{t}) + \mathbf{X}_{it}^T \gamma + \mu_i + \lambda_t + \epsilon_{it}
%\end{equation}
%Where \(Y_{it}\) represents the outcome variable for municipality \(i\) at time \(t\), measuring aspects of anti-competitive behavior such as the size of the jump in succession rate at the threshold.\(\text{Audit}_{it}\) is an indicator variable equal to 1 if municipality \(i\) was audited in year \(t\), and 0 otherwise.\(\text{Post}_{t}\) is a time dummy equal to 1 for all time periods following the implementation of the audit program, and 0 before.\(\delta (\text{Audit}_{it} \times \text{Post}_{t})\) is the interaction term between the audit and post variables; \(\delta\) measures the differential effect of being audited in the post-audit period, capturing the causal impact of the audits. \(\mathbf{X}_{it}\) is a vector of control variables for municipality \(i\) at time \(t\), such as economic factors, demographic characteristics, and other policies that might affect the outcome. \(\gamma\) is a vector of coefficients for the control variables. \(\mu_i\) and \(\lambda_t\) are fixed effects for municipalities and time, respectively, to control for invariant characteristics across municipalities and common temporal shocks. \(\epsilon_{it}\) is the error term.


\section{Main results (TBC)}\label{res}

\begin{figure}[htbp!]\centering\caption{Anti-competitive Behavior}\label{fig:anti-comp}
\subfloat[Incumbancy ]{\includegraphics[width=0.8\textwidth]{incumbent.png}}\\
\subfloat[Backlog 90 days ]{\includegraphics[width=0.8\textwidth]{logrol_std.png}}\\
\begin{minipage}{0.8\textwidth}
{\scriptsize \emph{Notes:} This figure displays the results from regression discontinuity design (RDD) analyses, with the margin of victory as the running variable, where negative values indicate a winning bid. In panel (a), the dependent variable is a dummy indicating whether the firm won the last auction it participated in. In panel (b), the dependent variable is the standardized backlog of the previous three months. \par}
\end{minipage}
\end{figure}


In this section, we show the results related to the potential presence of anticompetitive behavior in public procurement in the state of Sao Paulo in Brazil for the period 2009-2019. We start by graphically showing the aggregated results in Figure \ref{fig:anti-comp}. This figure displays the results from regression discontinuity design (RDD) analyses, using the \textit{margin of victory} as the running variable, where negative values indicate a winning bid. In panel (a), the dependent variable is a dummy indicating whether the firm won the last auction it participated in. In panel (b), the dependent variable is the standardized backlog of the previous three months. The results suggest that firms that narrowly won auctions exhibit a higher probability of winning the previous auction and have won more money in the past few months. Specifically, we find a coefficient of -0.171 (Std. Err. = 0.009) for the regression with incumbency as the dependent variable. Instead, when using the three months backlog, the coefficient was -0.271 (Std. Err. = 0.017).\footnote{The estimates are obtained using the \texttt{rdrobust} command in Stata, with standard errors clustered by auction.} This pattern of consecutive wins and considerable past earnings might suggest preferential treatment rather than equitable collaboration among companies. In collusion scenarios, losing firms would typically have higher earnings due to the shared nature of auction winnings among colluding parties. Instead, the results indicate that certain firms are consistently favored, likely due to political connections, which enables their repeated success in securing contracts and accumulating higher revenues. 





%\section{Audits effects (TBC)}\label{res}
%In this section, we show the potential effects of audit disclosure on anti-competitive behavior. 


\section{Additional results (TBC)}\label{res_ar}
\subsection{Heterogenous effects}
\subsection{PCC the role of organized crime}
\subsection{Correlation with predicted corruption}
\subsection{Firms outcomes}

\section{Conclusion (TBC)}\label{concl}
Our analysis of public procurement in the state of São Paulo, Brazil, from 2009 to 2019, reveals significant evidence of anticompetitive behavior. Using a Regression Discontinuity Design (RDD) approach, we observed that firms narrowly winning auctions are likelier to have won previous auctions and exhibit substantial earnings in recent months.

These results suggest preferential treatment rather than a fair, competitive process. The consistency of wins and higher earnings among certain firms likely points to political connections influencing the outcomes, deviating from the expected equitable distribution of opportunities in a competitive market. This scenario contrasts with typical collusive behaviors where auction winnings are shared among participating firms, leading to higher earnings for losing firms as well.


\clearpage
 \newpage
\bibliographystyle{chicago}
\bibliography{biblio}
\newpage


\begin{subappendices}
\setcounter{table}{0} 
\setcounter{section}{0}
\setcounter{equation}{0} 
\setcounter{figure}{0} 
\renewcommand{\thetable}{\Alph{section}.\arabic{table}}
\renewcommand{\thesection}{\Alph{section}}
\renewcommand{\thesubsection}{\Alph{section}.\arabic{subsection}}
\renewcommand{\thefigure}{\Alph{section}.\arabic{figure}}

\clearpage

\section{Appendix}\label{ApA}

\end{subappendices}
\end{document}
